\documentclass{article}
\usepackage[utf8]{inputenc}
%\usepackage[english]{babel}
\usepackage[portuguese]{babel}
\usepackage{graphicx}
\usepackage[hidelinks]{hyperref}
\usepackage{amsmath} %lets us use \begin{proof}
%\usepackage[]{amssymb} %gives us the character \varnothing
\usepackage[dvipsnames]{xcolor}
\usepackage{verbatim}

\title{Universidade Federal de Pernambuco \\
Finanças Internacionais - Lista 3}
\author{Professor: Marcelo E. A. Silva \\ Aluno: Paulo F. S. Junior}
\date{14 de outubro de 2022}

\begin{document}
\maketitle %This command prints the title based on information entered above

%Section and subsection automatically number unless you put the asterisk next to them.

\section*{Questão 1}

\subsection*{Derivações}
Dado que os países têm preferências e restrições orçamentárias idênticas (em estrutura/forma), portanto resolvemos para o país i, onde $i=\{U, E\}$.
\\\\
Preferências das famílias e restrições (sob livre mobilidade de capital) são:
   $$U^{i}=\ln C_1^i + \ln C_2^i$$
   $$C_1^{i}+B_1^{i}=Q_1^{i}$$
   $$C_2^{i}=(1+r^*)B_1^{i}+Q_1^{i}$$
   $$C_1^{i}+\frac{C_2^{i}}{(1+r^*)}=Q_1^{i}+\frac{Q_2^{i}}{(1+r^*)}$$
\\
Reescrevendo a R.O.I em termos de $C_2^{i}$ e substituindo na função utilidade temos o problema das famílias dado como se segue (maximização não condicionada):
\begin{equation}
    \ln C_1^i + \ln[(1+r^*)(Q_1^{i}-C_1^{i})+Q_2^{i}]
\end{equation}
Realizando a condição de primeira ordem para consumo do primeiro período:
\begin{equation}
    \frac{\partial U^{i}}{\partial C_1^{i}} = \frac{1}{C_1^{i}}-\frac{(1+r^*)}{[(1+r^*)(Q_1^{i}-C_1^{i})+Q_2^{i}]} = 0
\end{equation}
Resolvendo para $C_1^{i}:$
\begin{equation}
    \frac{1}{C_1^{i}}=\frac{(1+r^*)}{[(1+r^*)(Q_1^{i}-C_1^{i})+Q_2^{i}]}
\end{equation}
\begin{equation}
    (1+r^*)(Q_1^{i}-C_1^{i})+Q_2^{i}=C_1^{i}(1+r^*)
\end{equation}
\begin{equation}
    (1+r^*)Q_1^{i}+Q_2^{i}=C_1^{i}2(1+r^*)
\end{equation}
\begin{equation}
    C_1^{i}=\frac{1}{2}\left[Q_1^{i}+\frac{Q_2^{i}}{(1+r^*)}\right]
\end{equation}
Em uma economia de dotação a conta corrente do país i será dada por:
\begin{equation}
    CA_1^{i}=Q_1^{i}-C_1^{i}
\end{equation}
e portanto, no ótimo:
\begin{equation}\label{eq:CAi}
    CA_1^{i}=\frac{1}{2}\left[Q_1^{i}-\frac{Q_2^{i}}{(1+r^*)}\right]
\end{equation}
A conta corrente "mundial" é definida como:
\begin{equation}
    CA^U+CA^E=0
\end{equation}
Portanto, no ótimo em $t=1$ temos:
\begin{equation}
    \frac{1}{2}\left[Q_1^{U}-\frac{Q_2^{U}}{(1+r^*)}\right] + 
    \frac{1}{2}\left[Q_1^{E}-\frac{Q_2^{E}}{(1+r^*)}\right]=0
\end{equation}
Isolando $r^*$ encontramos a taxa de juros mundial de equilíbrio:
\begin{equation}
    \frac{1}{2}\left[Q_1^{U}-\frac{Q_2^{U}}{(1+r^*)}\right] + 
    \frac{1}{2}\left[Q_1^{E}-\frac{Q_2^{E}}{(1+r^*)}\right]=0 *(2)
\end{equation}
\begin{equation}
    Q_1^{U} + Q_1^{E}=\frac{Q_2^{E}+Q_2^{U}}{(1+r^*)}
\end{equation}
\begin{equation}
    (1+r^*)=\frac{Q_2^{E}+Q_2^{U}}{Q_1^{U} + Q_1^{E}}
\end{equation}
\begin{equation}\label{eq:r*equilibrio}
    r^*=\frac{Q_2^{E}+Q_2^{U}}{Q_1^{U} + Q_1^{E}}-1
\end{equation}
%Portanto, em equilíbrio, as contas correntes seriam dadas por:


\subsection*{Letra A)}
Com dotações $Q_1^U=Q_2^U=Q_1^E=Q_2^E=10$, e da derivação feita temos que:\\
A taxa de juros de equilíbrio \eqref{eq:r*equilibrio} será:
\begin{equation}
    r^*=\frac{10+10}{10+10}-1=0
\end{equation}
Utilizando a taxa de juros de equilíbrio encontrada acima, a conta corrente \eqref{eq:CAi} em equilíbrio dos países será de:
\begin{equation}
    CA_1^{E}=CA_1^{U}=\frac{1}{2}\left[Q_1^{i}-\frac{Q_2^{i}}{(1+r^*)}\right]
\end{equation}
\begin{equation}
    CA_1^{E}=CA_1^{U}=\frac{1}{2}\left[10-\frac{10}{(1+0)}\right]=0
\end{equation}
\subsection*{Letra B)}
Com a contração temporária nos EUA, $Q_1^{U}=8$, mas, $Q_2^U=Q_1^E=Q_2^E=10$,\\
A taxa de juros mundial de equilíbrio \eqref{eq:r*equilibrio} se ajustará para:
\begin{equation}
    r^*=\frac{10+10}{8 + 10}-1 = 0.11111111
\end{equation}
A conta corrente \eqref{eq:CAi} em equilíbrio dos EUA será de:
\begin{equation}
    CA_1^{U}=\frac{1}{2}\left[8-\frac{10}{(1+0,11111111)}\right] = \frac{1}{2}\left[8-9\right] = - 0.5
\end{equation}
A conta corrente \eqref{eq:CAi} em equilíbrio da Europa será de:
\begin{equation}
    CA_1^{E}=\frac{1}{2}\left[10-\frac{10}{(1+0,11111111)}\right] = \frac{1}{2}\left[10-9\right] = + 0.5
\end{equation}
Intuição:
A queda na dotação das famílias estado-unidenses no período um leva a antecipação do consumo (como forma de suavização), ou seja, trazem renda do futuro para o presente, o que portanto eleva a demanda por fundos no mercado mundial, o excesso de fundos então é eliminado através do aumento da taxa de juros mundial induzindo o mercado ao equilíbrio pois uma maior taxa de juros incetiva os europeus a pouparem mais hoje para consumir no futuro.

\subsection*{Letra C)}
Com a contração nos EUA, $Q_1^{U}=8$ e $Q_2^{U}=6$, mas, $Q_1^E=Q_2^E=10$,\\
A taxa de juros mundial de equilíbrio \eqref{eq:r*equilibrio} se ajustará para:
\begin{equation}
    r^*=\frac{10+6}{8 + 10}-1 = -0.11111111
\end{equation}
A conta corrente \eqref{eq:CAi} em equilíbrio dos EUA será de:
\begin{equation}
    CA_1^{U}=\frac{1}{2}\left[8-\frac{6}{(1-0,11111111)}\right] = \frac{1}{2}\left[8-6.75\right] = 0.625
\end{equation}
A conta corrente \eqref{eq:CAi} em equilíbrio da Europa será de:
\begin{equation}
    CA_1^{E}=\frac{1}{2}\left[10-\frac{10}{(1-0,11111111)}\right] = \frac{1}{2}\left[10-11.25\right] = -0.625
\end{equation}
Intuição:
Ne letra B a queda na dotação era temporária, embora caísse em um período, no outro a dotação se elevaria, por conseguinte a abundancia futura mundial seria maior que a corrente e portanto explica a elevação da taxa de juros mundial.
Na letra C, com a contração esperada no período dois sendo pior que o período um as famílias, estado-unidenses irão suavizar consumo poupando parte da sua dotação que portanto eleva a oferta de fundos no mercado mundial, o excesso de oferta portanto é eliminado através da redução de juros mundial (para um nível negativo visto que a dotação do segundo período é menor que a do primeiro), com a taxa de juros negativa os europeus têm incentivos para consumir mais hoje do que sua dotação corrente. 

\subsection*{Letra D)}
Com base no modelo acima podemos argumentar que as perspectivas futuras eram negativas e demoradas, ou seja, as dotações futuras seriam menores do que as dotações correntes durante algum período (relativamente longo) de tempo o que leva a queda da taxa de juros mundial, dessa forma os indivíduos buscavam liquidez.

\section*{Questão 2}
\subsection*{Letra A)}
Com dotações $Q_1^{U}=Q_1^{E}=8$ e $Q_2^U=Q_2^E=10$, teremos a taxa de juros mundial de equilíbrio \eqref{eq:r*equilibrio} de:
\begin{equation}
    r^*=\frac{10+10}{8 + 8}-1 = 0,25
\end{equation}
\subsection*{Letra B)}
Com dotações $Q_1^{U}=Q_1^{E}=Q_2^U=Q_2^E=8$, teremos a taxa de juros mundial de equilíbrio \eqref{eq:r*equilibrio} de:
\begin{equation}
    r^*=\frac{8+8}{8 + 8}-1 = 0
\end{equation}
\subsection*{Letra C)}
Com dotações $Q_1^{U}=Q_1^{E}=8$, mas, $Q_2^U=Q_2^E=6$, teremos a taxa de juros mundial de equilíbrio \eqref{eq:r*equilibrio} de:
\begin{equation}
    r^*=\frac{6+6}{8 + 8}-1 = -0,25
\end{equation}
\subsection*{Letra D)}
(Uma análise bem simples [serei demitido?!])\\
A equipe apresentou três cenários prováveis:
\begin{itemize}
    \item Cenário A\\
    Há uma contração do produto corrente porém a economia volta ao "normal" no período dois, com isso a taxa de juros é esperada subir.
    \item Cenário B\\
    Há uma contração do produto corrente e futuro de mesma proporção, dessa forma é esperado que a taxa de juros seja 0.
    \item Cenário C\\
    Há uma contração do produto nos dois períodos de tempo, porém, a queda do período futuro é maior do que a do período corrente, e portanto é esperada que a taxa de juros caia (para um nível negativo).
\end{itemize}
Pode-se levar em consideração os 3 casos e portando determinar uma possível expectativa com as três possibilidades possíveis apresentada pela equipe:
\begin{equation}
    E(r^*)=\frac{1}{3}\left[0,25+0-0,25\right]=0
\end{equation}
Entre pessimismo e otimismo, o cenário C, seria o cenário mais razoável.

\section*{Questão 3}
\subsection*{Ponto i)}
As famílias estado-unidenses tem preferências e restrições dadas por:
    $$U^{U}=\ln C_1^U + \beta\ln C_2^U$$
   $$C_1^{U}+B_1^{U}=Q_1^{U}$$
   $$C_2^{U}=(1+r_1^U)B_1^{U}+Q_1^{U}$$
   $$C_1^{U}+\frac{C_2^{U}}{(1+r_1^U)}=Q_1^{U}+\frac{Q_2^{U}}{(1+r_1^U)}$$
O problema da família portanto se resume em:
\begin{equation}
    Max L = \ln C_1^U + \beta\ln C_2^U-\lambda\left[C_1^{U}+\frac{C_2^{U}}{(1+r_1^U)}-Q_1^{U}-\frac{Q_2^{U}}{(1+r_1^U)}\right]
\end{equation}
Condições de primeira ordem:
\begin{equation}
    \frac{\partial L}{\partial C_1^U} = \frac{1}{C_1^U}-\lambda \to \lambda=\frac{1}{C_1^U}
\end{equation}
\begin{equation}
    \frac{\partial L}{\partial C_2^U} = \frac{\beta}{C_2^U}-\frac{\lambda}{(1+r_1^U)} \to \lambda=\frac{\beta(1+r_1^U)}{C_2^U}
\end{equation}
Para $\lambda=\lambda$ temos a equação de Euler:
\begin{equation}\label{eq:3aeuler}
    C_2^U=\beta(1+r_1^U)C_1^U
\end{equation}
Substituindo portanto \eqref{eq:3aeuler} na R.O.I, e isolando $C_1^U$, temos:
\begin{equation}
    C_1^{U}+\frac{\beta(1+r_1^U)C_1^U}{(1+r_1^U)}=Q_1^{U}+\frac{Q_2^{U}}{(1+r_1^U)}
\end{equation}
\begin{equation}
    (1+\beta)C_1^{U}=Q_1^{U}+\frac{Q_2^{U}}{(1+r_1^U)}
\end{equation}
Portanto, o consumo das famílias estado-unidense no primeiro período será:
\begin{equation}
    C_1^{U}=\left[Q_1^{U}+\frac{Q_2^{U}}{(1+r_1^U)}\right]\frac{1}{(1+\beta)}
\end{equation}
Utilizando a equação de Euler, \eqref{eq:3aeuler} teremos que $C_2^U$ será:
\begin{equation}
    C_2^{U}=\left[Q_1^{U}+\frac{Q_2^{U}}{(1+r_1^U)}\right]\frac{\beta(1+r_1^U)}{(1+\beta)}
\end{equation}
Da restrição no período um, podemos reescreve-la como:
\begin{equation}
    B_1^{U}=Q_1^{U}-C_1^{U}
\end{equation}
No ótimo:
\begin{equation}
    B_1^{U}=Q_1^{U}-\left[Q_1^{U}+\frac{Q_2^{U}}{(1+r_1^U)}\right]\frac{1}{(1+\beta)}
\end{equation}

\subsection*{Ponto ii)}
As famílias do resto do mundo tem preferências e restrições dadas por:
    $$U^{R}=\ln C_1^R + \beta\ln C_2^R$$
   $$C_1^{R}+B_1^{R}=Q_1^{R}$$
   $$C_2^{R}=(1+r_1^R)B_1^{R}+Q_1^{R}$$
   $$C_1^{R}+\frac{C_2^{R}}{(1+r_1^R)}=Q_1^{R}+\frac{Q_2^{R}}{(1+r_1^R)}$$
O problema da família portanto se resume em:
\begin{equation}
    Max L = \ln C_1^R + \beta\ln C_2^R-\lambda\left[C_1^{R}+\frac{C_2^{R}}{(1+r_1^R)}-Q_1^{R}-\frac{Q_2^{R}}{(1+r_1^R)}\right]
\end{equation}
Condições de primeira ordem:
\begin{equation}
    \frac{\partial L}{\partial C_1^R} = \frac{1}{C_1^R}-\lambda \to \lambda=\frac{1}{C_1^R}
\end{equation}
\begin{equation}
    \frac{\partial L}{\partial C_2^R} = \frac{\beta}{C_2^R}-\frac{\lambda}{(1+r_1^R)} \to \lambda=\frac{\beta(1+r_1^R)}{C_2^R}
\end{equation}
Para $\lambda=\lambda$ temos a equação de Euler:
\begin{equation}\label{eq:3beuler}
    C_2^R=\beta(1+r_1^R)C_1^R
\end{equation}
Substituindo portanto \eqref{eq:3beuler} na R.O.I, e isolando $C_1^R$, temos:
\begin{equation}
    C_1^{R}+\frac{\beta(1+r_1^R)C_1^R}{(1+r_1^R)}=Q_1^{R}+\frac{Q_2^{R}}{(1+r_1^R)}
\end{equation}
\begin{equation}
    (1+\beta)C_1^{R}=Q_1^{R}+\frac{Q_2^{R}}{(1+r_1^R)}
\end{equation}
Portanto, o consumo das famílias do resto do mundo no primeiro período será:
\begin{equation}
    C_1^{R}=\left[Q_1^{R}+\frac{Q_2^{R}}{(1+r_1^R)}\right]\frac{1}{(1+\beta)}
\end{equation}
Utilizando a equação de Euler, \eqref{eq:3beuler} teremos que $C_2^R$ será:
\begin{equation}
    C_2^{R}=\left[Q_1^{R}+\frac{Q_2^{R}}{(1+r_1^R)}\right]\frac{\beta(1+r_1^R)}{(1+\beta)}
\end{equation}
Da restrição no período um, podemos reescreve-la como:
\begin{equation}
    B_1^{R}=Q_1^{R}-C_1^{R}
\end{equation}
No ótimo:
\begin{equation}
    B_1^{R}=Q_1^{R}-\left[Q_1^{R}+\frac{Q_2^{R}}{(1+r_1^R)}\right]\frac{1}{(1+\beta)}
\end{equation}

\subsection*{Ponto iii)}
A livre mobilidade de capitais implica que $r_1^U=r*$. E, como não há posição líquida de ativos no período zero, $CA_1^U = TB_1^U = B_1^{U} = Q_1^U-C_1^U$, portanto:
\begin{equation}\label{eq:33ca1}
    CA_1^U = TB_1^U = B_1^{U} = Q_1^{U}-\left[Q_1^{U}+\frac{Q_2^{U}}{(1+r^*)}\right]\frac{1}{(1+\beta)}
\end{equation}
Para o período dois, teremos que $B_2^U=0$ pela TVC, além do mais:
\begin{equation}
    TB_2^U = Q_2^U - C_2^U = Q_2^U - \left[Q_1^{U}+\frac{Q_2^{U}}{(1+r_1^U)}\right]\frac{\beta(1+r_1^U)}{(1+\beta)}
\end{equation}
Para a conta corrente:
\begin{equation}
    CA_2^U = B_2^U - B_1^U = -B_1^U = \left[Q_1^{U}+\frac{Q_2^{U}}{(1+r^*)}\right]\frac{1}{(1+\beta)}-Q_1^{U}
\end{equation}

\subsection*{Ponto iv)}
A livre mobilidade de capitais implica que $r_1^R=r*$. Supondo que não há posição líquida de ativos no período zero, $CA_1^R = TB_1^R = B_1^{R} = Q_1^R-C_1^R$, portanto:
\begin{equation}
    CA_1^R = TB_1^R = B_1^{R} = Q_1^{R}-\left[Q_1^{R}+\frac{Q_2^{R}}{(1+r^*)}\right]\frac{1}{(1+\beta)}
\end{equation}
Para o período dois, teremos que $B_2^R=0$ pela TVC, além do mais:
\begin{equation}
    TB_2^R = Q_2^R - C_2^R = Q_2^R - \left[Q_1^{R}+\frac{Q_2^{R}}{(1+r_1^R)}\right]\frac{\beta(1+r_1^R)}{(1+\beta)}
\end{equation}
Para a conta corrente:
\begin{equation}
    CA_2^R = B_2^R - B_1^R = -B_1^R = \left[Q_1^{R}+\frac{Q_2^{R}}{(1+r^*)}\right]\frac{1}{(1+\beta)}-Q_1^{R}
\end{equation}

\subsection*{Ponto v)}
Sabemos que, "o recurso poupado por um país é a utilizado por outro país":
\begin{equation}
    CA_t^U+CA_t^R=0
\end{equation}
Portanto, para o período um, teremos:
\begin{equation}\label{eq:3vca1}
    Q_1^{U}-\left[Q_1^{U}+\frac{Q_2^{U}}{(1+r^*)}\right]\frac{1}{(1+\beta)} +
    Q_1^{R}-\left[Q_1^{R}+\frac{Q_2^{R}}{(1+r^*)}\right]\frac{1}{(1+\beta)}=0
\end{equation}
\begin{equation}
    Q_1^{U}+Q_1^{R}-\left[\frac{1}{1+\beta}\right] \left[Q_1^{U}+\frac{Q_2^{U}}{(1+r^*)}+Q_1^{R}+\frac{Q_2^{R}}{(1+r^*)}\right]=0
\end{equation}
Consequentemente, para o período dois, teremos:
\begin{equation}
    \left[Q_1^{U}+\frac{Q_2^{U}}{(1+r^*)}\right]\frac{1}{(1+\beta)}-Q_1^{U}+
    \left[Q_1^{R}+\frac{Q_2^{R}}{(1+r^*)}\right]\frac{1}{(1+\beta)}-Q_1^{R}=0
\end{equation}
\begin{equation}
    \left[\frac{1}{(1+\beta)}\right]\left[Q_1^{U}+\frac{Q_2^{U}}{(1+r^*)}+Q_1^{R}+\frac{Q_2^{R}}{(1+r^*)}\right]-Q_1^{U}-Q_1^{R}=0
\end{equation}

\subsection*{Ponto vi)}
Partindo da equação \eqref{eq:3vca1}, isolando $r^*$, teremos em equilíbrio:
\begin{equation}
    Q_1^{U}+Q_1^{R}-\left[\frac{1}{1+\beta}\right] \left[Q_1^{U}+\frac{Q_2^{U}}{(1+r^*)}+Q_1^{R}+\frac{Q_2^{R}}{(1+r^*)}\right]=0
\end{equation}
\begin{equation}
    \frac{(1+\beta)Q_1^U+(1+\beta)Q_1^R-Q_1^U-Q_1^R}{1+\beta}=
    \left[\frac{1}{(1+\beta)(1+r^*)}\right] 
    \left[Q_2^{U}+Q_2^{R}\right]
\end{equation}
\begin{equation}
    (1+r^*)=
    \frac{(1+\beta)}{\beta(Q_1^U+ Q_1^R)}
    \left[\frac{1}{(1+\beta)}\right] 
    \left[Q_2^{U}+Q_2^{R}\right]
\end{equation}
\begin{equation}
    (1+r^*)=
    \frac{Q_2^{U}+Q_2^{R}}{\beta(Q_1^U+Q_1^R)}
\end{equation}
Portanto a taxa de juros mundial de equilíbrio é de:
\begin{equation}\label{eq:3r*star}
    r^*=
    \frac{Q_2^{U}+Q_2^{R}}{\beta(Q_1^U+Q_1^R)}-1
\end{equation}
Utilizando a equação \eqref{eq:3r*star} acima e isolando $\beta(1+r^*)$:
\begin{equation}
    (1+r^*)=
    \frac{Q_2^{U}+Q_2^{R}}{\beta(Q_1^U+Q_1^R)}
\end{equation}
\begin{equation}
    \beta(1+r^*)=
    \frac{Q_2^{U}+Q_2^{R}}{Q_1^U+Q_1^R}
\end{equation}
Para que $\beta(1+r^*)=1$ as dotações mundiais ($Q_t^U+Q_t^R$) devem ser iguais.\\
Pelas equações de Euler \eqref{eq:3aeuler} e \eqref{eq:3beuler} isso implica que o consumo será igual em ambos os períodos e países.

\subsection*{Ponto vii)}
Partindo da taxa de juros de equilíbrio \eqref{eq:3r*star}, e substituindo o efeito em $Q_1^i$:
\begin{equation}
    r^*=    \frac{Q_2^{U}+Q_2^{R}}{\beta(1+x)(Q_1^U+Q_1^R)}-1
\end{equation}
O efeito marginal do aumento em x na taxa de juros mundial de equilíbrio será conforme:
\begin{equation}
    \frac{\partial r^*}{\partial x}=-\frac{Q_2^{U}+Q_2^{R}}{\beta(1+x)^2(Q_1^U+Q_1^R)}<0
\end{equation}
Um aumento da dotação mundial do primeiro período levará a uma queda na taxa de juros mundial de equilíbrio.
\begin{equation}
    CA_1^U = (1+x)Q_1^{U}-\left[(1+x)Q_1^{U}+\frac{Q_2^{U}}{(1+r^{*New})}\right]\frac{1}{(1+\beta)}
\end{equation}
\begin{equation}
    CA_1^U = Q_1^{U}-\left[Q_1^{U}+\frac{Q_2^{U}}{(1+r^*)}\right]\frac{1}{(1+\beta)}
\end{equation}
\begin{equation}
    \Delta CA_1^U = \left[x Q_1^{U}-\frac{x Q_1^{U}}{(1+\beta)}\right]+\left[\frac{Q_2^{U}}{(1+r^*)}-\frac{Q_2^{U}}{(1+r^{*New})}\right]\frac{1}{(1+\beta)}
\end{equation}
A primeira parte da variação é positiva enquanto a segunda, negativa, ou seja:
\begin{itemize}
    \item Caso 1:\\
    Se a variação da dotação corrente for maior do que a variação da dotação futura em valor presente, a conta corrente melhora, pois famílias dos EUA ofertam fundos.
    \item Caso 2:\\
    Se a variação da dotação corrente for igual a variação da dotação futura em valor presente, a conta corrente não se altera, ou seja, não há piora nem melhora.
    \item Caso 3:\\
    Se a variação da dotação corrente for menor a variação da dotação futura em valor presente, a conta corrente piora, pios famílias dos EUA demandam fundos.
\end{itemize}


\end{document}